\ifdefined\gkver\else\def\gkver{student}\fi
\documentclass[\gkver]{gaokaozhenti}
\begin{document}

\chapter{2019年海南理综}
%% paperType: 理综物理部分

\begin{enumerate}
\item
%% number: 1
%% paperName: 2019年海南理综第 1 题 4 分
%% typeId: 1
%% type: 单选题
%% chapter: 电场
%% point: 电场线
%% method: 无
%% score: 4
%% degree: 750
%% duplicateId: 0
%% body:
如图，静电场中的一条电场线上有 M、N 两点，箭头代表电场的方向，则\xzanswer{C}

\onepicture[w=4.5cm]{\includesvg{figs/01}}

\fourchoices[answer=C]
{M 点的电势比 N 点的低}
{M 点的场强大小一定比 N 点的大}
{电子在 M 点的电势能比在 N 点的低}
{电子在 M 点受到的电场力大小一定比在 N 点的大}

%% answer: C
%% memo:
\memoanswer{
沿着电场线方向电势逐渐降低，则 M 点电势比 N 点的高，故 A 错误；电场线的疏密程度反映电场强度的大小，图中只有一条电场线，所以无法判断 M 点和 N 点电场强度的大小，故 B 错误；电场力 $F=qE$，所以电场力的大小也无法判断，故 D 错误；电势能 $E_p=q\varphi$，负电荷在电势高的地方电势能小，故 C 正确。
}

\item
%% number: 2
%% paperName: 2019年海南理综第 2 题 4 分
%% typeId: 1
%% type: 单选题
%% chapter: 磁场
%% point: 左手定则
%% method: 等效替代
%% score: 4
%% degree: 650
%% duplicateId: 0
%% body:
如图，一段半圆形粗铜线固定在绝缘水平桌面（纸面）上，铜线所在空间有一匀强磁场，磁场方向竖直向下。当铜线通有顺时针方向电流时，铜线所受安培力的方向\xzanswer{A}

\onepicture[w=4.5cm]{figs/02.png}

\fourchoices[answer=A]
{向前}
{向后}
{向左}
{向右}

%% answer: A
%% memo:
\memoanswer{
由左手定则和力的叠加可知，安培力方向向前，故 A 正确。
}

\item
%% number: 3
%% paperName: 2019年海南理综第 3 题 4 分
%% typeId: 1
%% type: 单选题
%% chapter: 直线运动
%% point: 匀变速直线运动
%% method: 无
%% score: 4
%% degree: 750
%% duplicateId: 0
%% body:
汽车在平直公路上以 $20\Ums$ 的速度匀速行驶。前方突遇险情，司机紧急刹车，汽车做匀减速运动，加速度大小为 $8\Umsq$。从开始刹车到汽车停止，汽车运动的距离为\xzanswer{C}

\fourchoices[answer=C]
{$10\Um$}
{$20\Um$}
{$25\Um$}
{$50\Um$}

%% answer: C
%% memo:
\memoanswer{
由题意可知，刹车时汽车的初速度 $v_0=20\Ums$，末速度 $v=0$，加速度为 $a=-8\Umsq$，由速度位移公式得 $v^2-v_0^2=2ax$，解得 $x=25\Um$，故 C 正确。
}

\item
%% number: 4
%% paperName: 2019年海南理综第 4 题 4 分
%% typeId: 1
%% type: 单选题
%% chapter: 曲线运动
%% point: 人造地球卫星
%% method: 无
%% score: 4
%% degree: 650
%% duplicateId: 0
%% body:
2019 年 5 月，我国第 45 颗北斗卫星发射成功。已知该卫星轨道距地面的高度约为 $36\,000\Ukm$，是“天宫二号”空间实验室轨道高度的 90 倍左右，则\xzanswer{B}

\fourchoices[answer=B]
{该卫星的速率比“天宫二号”的大}
{该卫星的周期比“天宫二号”的大}
{该卫星的角速度比“天宫二号”的大}
{该卫星的向心加速度比“天宫二号”的大}

%% answer: B
%% memo:
\memoanswer{
卫星绕地球做匀速圆周运动，设其运动轨道半径为 $r$，由万有引力提供向心力得 $G\frac{Mm}{r^2}=m\frac{v^2}{r}$，则 $v=\sqrt{\frac{GM}{r}}$，该卫星的轨道半径 $r$ 大，所以该卫星的速率比“天宫二号”的小，故 A 错误；由 $G\frac{Mm}{r^2}=m\left(\frac{2\pi}{T}\right)^2r$ 得 $T=2\pi\sqrt{\frac{r^3}{GM}}$，该卫星的轨道半径 $r$ 大，则该卫星的周期比“天宫二号”的大，故 B 正确；由 $G\frac{Mm}{r^2}=m\omega^2r$ 得 $\omega=\sqrt{\frac{GM}{r^3}}$，该卫星的轨道半径 $r$ 大，则该卫星的角速度比“天宫二号”的小，故 C 错误；由 $G\frac{Mm}{r^2}=ma$ 得 $a=\frac{GM}{r^2}$，该卫星的轨道半径 $r$ 大，则该卫星的加速度比“天宫二号”的小，故 D 错误。
}

\item
%% number: 5
%% paperName: 2019年海南理综第 5 题 4 分
%% typeId: 1
%% type: 单选题
%% chapter: 运动和力
%% point: 连接体
%% method: 无
%% score: 4
%% degree: 550
%% duplicateId: 0
%% body:
如图，两物块 P、Q 置于水平地面上，其质量分别为 $m$、$2m$，两者之间用水平轻绳连接。两物块与地面之间的动摩擦因数均为 $\mu$，重力加速度大小为 $g$，现对 Q 施加一水平向右的拉力 $F$，使两物块做匀加速直线运动，轻绳的张力大小为\xzanswer{D}

\onepicture[w=5.0cm]{\includesvg{figs/05}}

\fourchoices[answer=D]
{$F-2\mu mg$}
{$\frac{1}{3}F+\mu mg$}
{$\frac{1}{3}F-\mu mg$}
{$\frac{1}{3}F$}

%% answer: D
%% memo:
\memoanswer{
对 P、Q 整体受力分析，由牛顿第二定律得 $F-\mu mg-2\mu mg=3ma$，对物块 P 受力分析，根据牛顿第二定律有 $F_T-\mu mg=ma$，联立解得 $F_T=\frac{1}{3}F$，故 D 正确。
}

\item
%% number: 6
%% paperName: 2019年海南理综第 6 题 4 分
%% typeId: 1
%% type: 单选题
%% chapter: 曲线运动
%% point: 圆周运动的应用
%% method: 无
%% score: 4
%% degree: 550
%% duplicateId: 0
%% body:
如图，一硬币（可视为质点）置于水平圆盘上，硬币与竖直转轴 $OO'$ 的距离为 $r$，已知硬币与圆盘之间的动摩擦因数为 $\mu$（最大静摩擦力等于滑动摩擦力），重力加速度大小为 $g$。若硬币与圆盘一起绕 $OO'$ 轴匀速转动，则圆盘转动的最大角速度为\xzanswer{B}

\onepicture[w=4.5cm]{figs/06.png}

\fourchoices[answer=B]
{$\frac{1}{2}\sqrt{\frac{\mu g}{r}}$}
{$\sqrt{\frac{\mu g}{r}}$}
{$\sqrt{\frac{2\mu g}{r}}$}
{$2\sqrt{\frac{\mu g}{r}}$}

%% answer: B
%% memo:
\memoanswer{
对硬币受力分析可知，由静摩擦力提供向心力，达到最大角速度时，静摩擦力等于滑动摩擦力，则有 $\mu mg=m\omega^2r$，解得 $\omega=\sqrt{\frac{\mu g}{r}}$，故 B 正确。
}

\item
%% number: 7
%% paperName: 2019年海南理综第 7 题 5 分
%% typeId: 2
%% type: 多选题
%% chapter: 光学
%% point: 光电效应
%% method: 无
%% score: 5
%% degree: 550
%% duplicateId: 0
%% body:
对于钠和钙两种金属，其遏止电压 $U_c$ 与入射光频率 $\nu$ 的关系如图所示。用 $h$、$e$ 分别表示普朗克常量和电子电荷量，则\xzanswer{AB}

\onepicture[w=4.5cm]{\includesvg{figs/07}}

\fourchoices[answer=AB]
{钠的逸出功小于钙的逸出功}
{图中直线的斜率为 $\frac{h}{e}$}
{在得到这两条直线时，必须保证入射光的光强相同}
{若这两种金属产生的光电子具有相同的最大初动能，则照射到钠的光频率较高}

%% answer: AB
%% memo:
\memoanswer{
由光电效应方程有 $U_ce=h\nu-W_0$，变形得 $U_c=-W_0+\frac{h}{e}\nu$，图象与纵坐标的截距大小为 $\frac{W_0}{e}$，而 $e$ 为电子电荷量，所以钠的逸出功小于钙的逸出功，斜率为 $\frac{h}{e}$，故 AB 正确；遏止电压与光照强度没有关系，故 C 错误；由于 $E_{\text{km}}=h\nu-W_0$，若这两种金属产生光电子的最大初动能相同，则照射钙的光频率较高，故 D 错误。
}

\item
%% number: 8
%% paperName: 2019年海南理综第 8 题 5 分
%% typeId: 2
%% type: 多选题
%% chapter: 交流电
%% point: 变压器
%% method: 无
%% score: 5
%% degree: 550
%% duplicateId: 0
%% body:
如图，一理想变压器输入端接交流恒压源，输出端电路由 $R_1$、$R_2$ 和 $R_3$ 三个电阻构成。将该变压器原、副线圈的匝数比由 $5:1$ 改为 $10:1$ 后\xzanswer{CD}

\onepicture[w=5.0cm]{figs/08.png}

\fourchoices[answer=CD]
{流经 $R_1$ 的电流减小到原来的 $\frac{1}{4}$}
{$R_2$ 两端的电压增加到原来的 2 倍}
{$R_3$ 两端的电压减小到原来的 $\frac{1}{2}$}
{电阻上总的热功率减小到原来的 $\frac{1}{4}$}

%% answer: CD
%% memo:
\memoanswer{
由理想变压器中的关系 $\frac{U_1}{U_2}=\frac{n_1}{n_2}$ 知，匝数比由 $5:1$ 变成 $10:1$ 后，副线圈两端电压变成原来的 $\frac{1}{2}$，则流经 $R_1$ 的电流变为原来的 $\frac{1}{2}$，故 A 错误；流经 $R_2$、$R_3$ 的电流变为原来的 $\frac{1}{2}$，$R_2$、$R_3$ 的电压变为原来的 $\frac{1}{2}$，故 B 错误，C 正确；电阻上总的热功率 $P=I^2R$ 变为原来的 $\frac{1}{4}$，故 D 正确。
}

\item
%% number: 9
%% paperName: 2019年海南理综第 9 题 5 分
%% typeId: 2
%% type: 多选题
%% chapter: 磁场
%% point: 带电粒子在磁场中的运动
%% method: 无
%% score: 5
%% degree: 550
%% duplicateId: 0
%% body:
如图，虚线 MN 的右侧有方向垂直于纸面向里的匀强磁场；两电荷量相同的粒子 P、Q 从磁场边界的 M 点先后射入磁场，在纸面内运动。射入磁场时，P 的速度 $v_P$ 垂直于磁场边界，Q 的速度 $v_Q$ 与磁场边界的夹角为 $45^{\circ}$。已知两粒子均从 N 点射出磁场，且在磁场中运动的时间相同，则\xzanswer{AC}

\onepicture[w=5.0cm]{figs/09.png}

\fourchoices[answer=AC]
{P 和 Q 的质量之比为 $1:2$}
{P 和 Q 的质量之比为 $\sqrt{2}:1$}
{P 和 Q 速度大小之比为 $\sqrt{2}:1$}
{P 和 Q 速度大小之比为 $2:1$}

%% answer: AC
%% memo:
\memoanswer{
设 MN 距离为 $d$，则粒子 P 的运动半径为 $R_1=\frac{1}{2}d$，粒子 Q 的运动半径为 $R_2=\frac{\sqrt{2}}{2}d$，运动时间相同，则有 $\frac{\pi R_1}{v_P}=\frac{\pi R_2}{2v_Q}$，速度之比 $\frac{v_P}{v_Q}=\frac{\sqrt{2}}{1}$，故 C 正确 D 错误；由洛伦兹力提供向心力有 $qvB=m\frac{v^2}{R}$，$R=\frac{mv}{qB}$，则质量比 $\frac{m_P}{m_Q}=\frac{1}{2}$，故 A 正确 B 错误。
}

\item
%% number: 10
%% paperName: 2019年海南理综第 10 题 5 分
%% typeId: 2
%% type: 多选题
%% chapter: 曲线运动
%% point: 平抛运动
%% method: 极值
%% score: 5
%% degree: 450
%% duplicateId: 0
%% body:
三个小物块分别从 3 条不同光滑轨道的上端由静止开始滑下。已知轨道 1、轨道 2、轨道 3 的上端距水平地面的高度均为 $4h_0$；它们的下端水平，距地面的高度分别为 $h_1=h_0$、$h_2=2h_0$、$h_3=3h_0$，如图所示。若沿轨道 1、2、3 下滑的小物块的落地点到轨道下端的水平距离分别记为 $s_1$、$s_2$、$s_3$，则\xzanswer{BC}

\onepicture[w=6.0cm]{figs/10.png}

\fourchoices[answer=BC]
{$s_1>s_2$}
{$s_2>s_3$}
{$s_1=s_3$}
{$s_2=s_3$}

%% answer: BC
%% memo:
\memoanswer{
物块沿轨道 1 下滑时，由动能定理得 $3mgh_0=\frac{1}{2}mv_1^2$，从轨道 1 末端抛出后做平抛运动，竖直方向 $h_0=\frac{1}{2}gt^2$，则水平位移 $s_1=v_1t$，解得 $s_1=2\sqrt{3}h_0$，同理可解得 $s_2=4h_0$，$s_3=2\sqrt{3}h_0$，故 BC 正确。
}

\item
%% number: 11
%% paperName: 2019年海南理综第 11 题 8 分
%% typeId: 4
%% type: 实验题
%% chapter: 电路
%% point: 伏安法测电阻
%% method: 无
%% score: 8
%% degree: 550
%% duplicateId: 0
%% body:
用实验室提供的器材设计一个测量电流表内阻的电路。实验室提供的器材为：待测电流表 A（量程 $10\UmA$，内阻约为 $50\UO$），滑动变阻器 $R_1$，电阻箱 $R$，电源 $E$（电动势约为 $6\UV$，内阻可忽略），开关 $S_1$ 和 $S_2$，导线若干。

\twopicture[labelA=2019海南理综11a, labelB=2019海南理综11b, w=4.5cm]
{figs/11a.png}
{figs/11b.png}

\begin{enumerate}
	\item
	根据实验室提供的器材，在图（a）所示虚线框内将电路原理图补充完整，要求滑动变阻器起限流作用；
	\item
	将图（b）中的实物按设计的原理图连线；
	\item
	若实验提供的滑动变阻器有两种规格：

	① $10\UO$，额定电流 $2\UA$

	② $1500\UO$，额定电流 $0.5\UA$

	实验中应该取\tkanswer[1.5cm]{②}。（填“①”或“②”）
\end{enumerate}

%% answer:
\jdanswer{
\begin{enumerate}
	\item
	如图甲所示

	\item
	如图乙所示

	\item
	②

\end{enumerate}
}
%% memo:
\memoanswer{
（1）要求滑动变阻器起限流作用，则滑动变阻器串联在电路中。本实验的原理是用半偏法测电流表内阻，闭合开关 $S_1$，调节滑动变阻器滑片，使电流表满偏，然后闭合开关 $S_2$，调节电阻箱阻值，使电流表半偏，此时电阻箱读数即为电流表内阻。故电路图如答图甲所示。

\onepicture[w=4.5cm]{figs/11_an1.png}

（2）按照电路图依据电流走向连接实物图，如答图乙所示。

\onepicture[w=5.5cm]{figs/11_an2.png}

（3）本实验要求电流表满偏，根据 $E=I(R_A+R_1)$ 可得 $R_1$ 约为 $550\UO$，所以滑动变阻器应选②。
}

\item
%% number: 12
%% paperName: 2019年海南理综第 12 题 10 分
%% typeId: 4
%% type: 实验题
%% chapter: 力物体的平衡
%% point: 胡克定律
%% method: 无
%% score: 10
%% degree: 550
%% duplicateId: 0
%% body:
某同学利用图（a）的装置测量轻弹簧的劲度系数。图中，光滑的细杆和直尺水平固定在铁架台上，一轻弹簧穿在细杆上，其左端固定，右端与细绳连接；细绳跨过光滑定滑轮，其下端可以悬挂砝码（实验中，每个砝码的质量均为 $m=50.0\Ug$）。弹簧右端连有一竖直指针，其位置可在直尺上读出。实验步骤如下：

\onepicture[w=5.0cm, align=right]{figs/12a.png}

①在绳下端挂上一个砝码，调整滑轮，使弹簧与滑轮间的细线水平且弹簧与细杆没有接触；

②系统静止后，记录砝码的个数及指针的位置；

③逐次增加砝码个数，并重复步骤②（保持弹簧在弹性限度内）；

④用 $n$ 表示砝码的个数，$l$ 表示相应的指针位置，将获得的数据记录在表格内。

回答下列问题：

\begin{enumerate}
	\item
	根据下表的实验数据在图（b）中补齐数据点并作出 $l-n$ 图像。

	\begin{table}[htbp]
	\centering
	\begin{tblr}{|c|c|c|c|c|c|}
	\hline
	$n$ & 1 & 2 & 3 & 4 & 5 \\
	\hline
	$l/\Ucm$ & 10.48 & 10.96 & 11.45 & 11.95 & 12.40 \\
	\hline
	\end{tblr}
	\end{table}

	\onepicture[w=5.5cm, align=right]{\includesvg{figs/12b}}

	\item
	弹簧的劲度系数 $k$ 可用砝码质量 $m$、重力加速度大小 $g$ 及 $l-n$ 图线的斜率 $\alpha$ 表示，表达式为 $k=$\tkanswer[2cm]{$\frac{mg}{\alpha}$}。若 $g$ 取 $9.80\Umsq$，则本实验中 $k=$\tkanswer[1.5cm]{99.7}$\UNm$（结果保留 3 位有效数字）。
\end{enumerate}

%% answer:
\jdanswer{
\begin{enumerate}
	\item
	如图所示

	\item
	$\frac{mg}{\alpha}$，99.7（98.0～102）

\end{enumerate}
}
%% memo:
\memoanswer{
（1）在 $l-n$ 坐标系中补齐（2,10.96）、（5,12.40）两点，然后用平直的直线连接，使大多数的点都在直线上或均匀分布在线的两侧。

\onepicture[w=5.5cm]{\includesvg{figs/12_an}}

（2）对弹簧受力分析有 $nmg=k(l-l_0)$，$l_0$ 为弹簧原长，变形有 $l=l_0+\frac{mg}{k}n$，所以图象的斜率表示 $\frac{mg}{k}$，即 $\alpha=\frac{mg}{k}$，得 $k=\frac{mg}{\alpha}$，图象斜率 $\alpha=\frac{0.0295}{6}$，代入 $k=\frac{mg}{\alpha}$，解得 $k=99.7\UNm$。
}

\item
%% number: 13
%% paperName: 2019年海南理综第 13 题 10 分
%% typeId: 6
%% type: 计算题
%% chapter: 运动和力
%% point: 动量和能量
%% method: 无
%% score: 10
%% degree: 550
%% duplicateId: 0
%% body:
如图，用不可伸长轻绳将物块 a 悬挂在 O 点；初始时，轻绳处于水平拉直状态。现将 a 由静止释放，当物块 a 下摆至最低点时，恰好与静止在水平面上的物块 b 发生弹性碰撞（碰撞时间极短），碰撞后 b 滑行的最大距离为 $s$。已知 b 的质量是 a 的 3 倍，b 与水平面间的动摩擦因数为 $\mu$，重力加速度大小为 $g$。求

\onepicture[w=5.5cm, align=right]{\includesvg{figs/13}}

\begin{enumerate}
	\item
	碰撞后瞬间物块 b 速度的大小；
	\item
	轻绳的长度。
\end{enumerate}

%% answer:
\jdanswer{
\begin{enumerate}
	\item
	$\sqrt{2\mu gs}$

	\item
	$4\mu s$

\end{enumerate}
}
%% memo:
\memoanswer{
（1）设碰撞后物块 b 的速度为 $v_b$，由牛顿第二定律得 $\mu m_bg=m_ba$，即物块 b 在水平面上运动的加速度大小为 $a=\mu g$，物块 b 做匀减速直线运动，由运动学公式得 $v_b^2=2as$，联立解得 $v_b=\sqrt{2\mu gs}$。

（2）设轻绳长度为 $l$，碰撞前物块 a 速度为 $v$，质量为 $m$，碰后速度为 $v_a$，由动能定理得 $mgl=\frac{1}{2}mv^2$，两物块发生弹性碰撞，由动量守恒得 $mv=mv_a+3mv_b$，由机械能守恒得 $\frac{1}{2}mv^2=\frac{1}{2}mv_a^2+\frac{1}{2}\cdot 3mv_b^2$，联立解得 $l=4\mu s$。
}

\item
%% number: 14
%% paperName: 2019年海南理综第 14 题 16 分
%% typeId: 6
%% type: 计算题
%% chapter: 电磁感应
%% point: 双杆问题
%% method: 无
%% score: 16
%% degree: 450
%% duplicateId: 0
%% body:
如图，一水平面内固定有两根平行的长直金属导轨，导轨间距为 $l$；两根相同的导体棒 AB、CD 置于导轨上并与导轨垂直，长度均为 $l$；棒与导轨间的动摩擦因数为 $\mu$（最大静摩擦力等于滑动摩擦力）；整个装置处于匀强磁场中，磁感应强度大小为 $B$，方向竖直向下。从 $t=0$ 时开始，对 AB 棒施加一外力，使 AB 棒从静止开始向右做匀加速运动，直到 $t=t_1$ 时刻撤去外力，此时棒中的感应电流为 $i_1$；已知 CD 棒在 $t=t_0$（$0<t_0<t_1$）时刻开始运动，运动过程中两棒均与导轨接触良好。两棒的质量均为 $m$，电阻均为 $R$，导轨的电阻不计。重力加速度大小为 $g$。

\onepicture[w=8.0cm, align=right]{figs/14.png}

\begin{enumerate}
	\item
	求 AB 棒做匀加速运动的加速度大小；
	\item
	求撤去外力时 CD 棒的速度大小；
	\item
	撤去外力后，CD 棒在 $t=t_2$ 时刻静止，求此时 AB 棒的速度大小。
\end{enumerate}

%% answer:
\jdanswer{
\begin{enumerate}
	\item
	$\frac{2\mu mgR}{B^2l^2t_0}$

	\item
	$\frac{2\mu mgRt_1}{B^2l^2t_0}-\frac{2Ri_1}{Bl}$

	\item
	$\frac{4\mu mgRt_1}{B^2l^2t_0}-\frac{2Ri_1}{Bl}-2\mu g(t_2-t_1)$

\end{enumerate}
}
%% memo:
\memoanswer{
（1）CD 棒在 $t=t_0$ 时刻开始运动，说明 $t_0$ 时刻 CD 棒安培力等于最大静摩擦力，即 $BIl=\mu mg$；$t_0$ 时刻，AB 棒的速度为 $v=at_0$，AB 棒切割产生的电动势为 $E=Blv$，回路中的电流为 $I=\frac{E}{2R}$，联立解得 $a=\frac{2\mu mgR}{B^2l^2t_0}$。

（2）$t=t_1$ 时刻撤去外力，$t_1$ 时刻 AB 棒的速度为 $v_1=at_1$，产生的感应电动势为 $E_1=Blv_1$；设 $t_1$ 时刻 CD 棒的速度为 $v_2$，产生的感应电动势为 $E_2=Blv_2$，回路中的感应电流为 $i_1=\frac{E_1-E_2}{2R}$，联立解得 $v_2=\frac{2\mu mgRt_1}{B^2l^2t_0}-\frac{2Ri_1}{Bl}$。

（3）对 AB、CD 整体，以向右为正方向，设 CD 棒停止运动时，AB 棒速度为 $v_3$。由于 AB、CD 运动方向相同，通过电流方向相反，二者所受的安培力大小相等、方向相反，对整体应用动量定理时相互抵消，由动量定理得 $-2\mu mg(t_2-t_1)=mv_3-(mv_1+mv_2)$，即 $v_3=v_1+v_2-2\mu g(t_2-t_1)$，代入数据有 $v_3=\frac{4\mu mgRt_1}{B^2l^2t_0}-\frac{2Ri_1}{Bl}-2\mu g(t_2-t_1)$。
}

\item
%% number: 15
%% paperName: 2019年海南理综第 15 题 8 分
%% typeId: 6
%% type: 计算题
%% chapter: 气体性质
%% point: 气体连接体
%% method: 无
%% score: 8
%% degree: 550
%% duplicateId: 0
%% body:
[选修 3-3]（12 分）

\begin{enumerate}
	\item
	（4 分）一定量的理想气体从状态 M 出发，经状态 N、P、Q 回到状态 M，完成一个循环。从 M 到 N、从 P 到 Q 是等温过程；从 N 到 P、从 Q 到 M 是等容过程；其体积$-$温度图像（$V-T$ 图）如图所示。下列说法正确的是\xzanswer{BCE}

	\onepicture[w=4.5cm, align=right]{figs/15a.png}

	\fivechoices[answer=BCE]
	{从 M 到 N 是吸热过程}
	{从 N 到 P 是吸热过程}
	{从 P 到 Q 气体对外界做功}
	{从 Q 到 M 是气体对外界做功}
	{从 Q 到 M 气体的内能减少}
	\item
	（8 分）如图，一封闭的圆柱形容器竖直放置在水平地面上，一重量不可忽略的光滑活塞将容器内的理想气体分为 A、B 两部分，A 体积为 $V_A=4.0\times10^{-3}\Umc$，压强为 $p_A=47\UcmHg$；B 体积为 $V_B=6.0\times10^{-3}\Umc$，压强为 $p_B=52\UcmHg$。现将容器缓慢转至水平，气体温度保持不变，求此时 A、B 两部分气体的体积。

	\onepicture[w=6.0cm, align=right]{figs/15b.png}
\end{enumerate}

%% answer:
\jdanswer{
\begin{enumerate}
	\item
	BCE

	\item
	$3.76\times10^{-3}\Umc$，$6.24\times10^{-3}\Umc$

\end{enumerate}
}
%% memo:
\memoanswer{
（1）从 M 到 N 气体做等温变化，内能不变，体积减小，外界对气体做功，由热力学第一定律可知，该过程是气体的放热过程，故 A 错误；从 N 到 P 气体做等容变化，温度升高，内能增大，由热力学第一定律可知，该过程是气体的吸热过程，故 B 正确；从 P 到 Q 气体体积增大，气体对外界做功，故 C 正确；从 Q 到 M 气体做等容变化，气体不对外界做功，但温度降低，内能减小，故 D 错误 E 正确。

（2）初状态密闭气体 A，$p_A=47\UcmHg$，$V_A=4.0\times10^{-3}\Umc$；初状态密闭气体 B，$p_B=52\UcmHg$，$V_B=6.0\times10^{-3}\Umc$。容器水平后，密闭气体 A 和密闭气体 B 压强相等，对密闭气体 A，由玻意耳定律得 $p_AV_A=pV_A'$；对密闭气体 B，由玻意耳定律得 $p_BV_B=pV_B'$；$V_A+V_B=V_A'+V_B'$，联立解得 $V_A'=3.76\times10^{-3}\Umc$，$V_B'=6.24\times10^{-3}\Umc$。
}

\item
%% number: 16
%% paperName: 2019年海南理综第 16 题 8 分
%% typeId: 6
%% type: 计算题
%% chapter: 光学
%% point: 全反射
%% method: 无
%% score: 8
%% degree: 450
%% duplicateId: 0
%% body:
[选修 3-4]（12 分）

\begin{enumerate}
	\item
	（4 分）一列简谐横波沿 $x$ 轴正方向传播，周期为 $0.2\Us$，$t=0$ 时的波形图如图所示。下列说法正确的是\xzanswer{ABC}

	\onepicture[w=6.5cm, align=right]{figs/16a.png}

	\fivechoices[answer=ABC]
	{平衡位置在 $x=1.0\Um$ 处的质点的振幅为 $0.03\Um$}
	{该波的波速为 $10\Ums$}
	{$t=0.3\Us$ 时，平衡位置在 $x=0.5\Um$ 处的质点向 $y$ 轴正向运动}
	{$t=0.4\Us$ 时，平衡位置在 $x=-1.0\Um$ 处的质点处于波谷位置}
	{$t=0.5\Us$ 时，平衡位置在 $x=1.0\Um$ 处的质点加速度为零}
	\item
	（8 分）一透明材料制成的圆柱体的上底面中央有一球形凹陷，凹面与圆柱体下底面可透光，表面其余部分均涂有遮光材料。过圆柱体对称轴线的截面如图所示。O 点是球形凹陷的球心，半径 OA 与 OG 夹角 $\theta=120^{\circ}$。平行光沿轴线方向向下入射时，从凹面边缘 A 点入射的光线经折射后，恰好由下底面上 C 点射出。已知 AB=FG=$1\Ucm$，BC=$\sqrt{3}\Ucm$，OA=$2\Ucm$。

	\onepicture[w=5.0cm, align=right]{figs/16b.png}

	\begin{enumerate}
		\item
		求此透明材料的折射率；
		\item
		撤去平行光，将一点光源置于球心 O 点处，求下底面上有光出射的圆形区域的半径（不考虑侧面的反射光及多次反射的影响）。
	\end{enumerate}
\end{enumerate}

%% answer:
\jdanswer{
\begin{enumerate}
	\item
	ABC

	\item
	$\sqrt{3}$，$\frac{\sqrt{2}+\sqrt{6}}{2}\Ucm$

\end{enumerate}
}
%% memo:
\memoanswer{
（1）由波动图象可知，$x=1.0\Um$ 处质点的振幅为 $3\Ucm=0.03\Um$，故 A 正确；波长 $\lambda=2\Um$，周期 $T=0.2\Us$，波速 $v=\frac{\lambda}{T}=10\Ums$，故 B 正确；$t=0.3\Us$ 时，即 $t=1\frac{1}{2}T$，$x=0.5\Um$ 处的质点处于平衡位置并向 $y$ 轴正向运动，故 C 正确；$t=0.4\Us$ 时，即 $t=2T$，平衡位置在 $x=-1.0\Um$ 处的质点处于波峰位置，故 D 错误；$t=0.5\Us$ 时，即 $t=2\frac{1}{2}T$，平衡位置在 $x=1.0\Um$ 处的质点处于波谷，加速度最大，故 E 错误。

（2）（i）光路图如图甲所示，由几何关系可知，入射角 $\alpha=60^{\circ}$，折射角 $\beta=30^{\circ}$，由折射定律得 $n=\frac{\sin\alpha}{\sin\beta}$，解得 $n=\sqrt{3}$。

\onepicture[w=5.0cm]{figs/16_an1.jpg}

（ii）光路图如图乙所示，设射到 H 点的光线恰好发生全反射时，有 $\sin C=\frac{1}{n}$，可以求得 $\tan C=\frac{\sqrt{2}}{2}$，由几何关系得 $R=OE\cdot\tan C$，又 $OE=(1+\sqrt{3})\Ucm$，解得 $R=\frac{\sqrt{2}+\sqrt{6}}{2}\Ucm$。

\onepicture[w=5.0cm]{figs/16_an2.jpg}
}

\end{enumerate}

\end{document}
